\documentclass[11pt,a4paper,twoside,openany]{book}
\usepackage{logique-fancy}
\begin{document}
\begin{titlepage}
\thispagestyle{empty}
\begin{tikzpicture}[remember picture,overlay]
 \fill[LFnavy] (current page.north west) rectangle ([yshift=-11.0cm]current page.north east);
 \fill[LFgold] ([yshift=-11.1cm]current page.north west) rectangle ([yshift=-11.0cm]current page.north east);
 \draw[LFteal!40,line width=1.5pt] ([xshift=18mm,yshift=-21mm]current page.north west) rectangle ([xshift=-18mm,yshift=21mm]current page.south east);
 \draw[LFgold!85,line width=.7pt] ([xshift=21mm,yshift=-24mm]current page.north west) rectangle ([xshift=-21mm,yshift=24mm]current page.south east);
\end{tikzpicture}
\vspace*{1.1cm}
\begin{center}
{\large\sffamily\color{white} LICENCE 2 MATHÉMATIQUES}\par\vspace{1.1cm}
{\fontsize{34}{41}\selectfont\bfseries\sffamily\color{white} LOGIQUE\\[3mm] MATHÉMATIQUE}\par\vspace{.7cm}
{\large\color{white} Du langage à la démonstration}\par
\end{center}
\vspace{3.4cm}
\begin{center}
{\Large\bfseries Abdelhalim Azzouz}\par\vspace{6mm}
{\sffamily Cours magistral et travaux dirigés}\par\vspace{1.4cm}
\begin{tcolorbox}[width=.7\textwidth,colback=LFpaper,colframe=LFgold,boxrule=.75pt,arc=2mm]
\centering\sffamily\bfseries Polycopié\\[2mm]
Cours et travaux dirigés
\end{tcolorbox}
\vfill
{\sffamily\color{LFnavy} Année universitaire 2026--2027}
\end{center}
\end{titlepage}
\thispagestyle{empty}
\frontmatter
\chapter*{Avant-propos}
\addcontentsline{toc}{chapter}{Avant-propos}
La logique mathématique n'est pas seulement un répertoire de symboles : elle fournit des règles permettant de formuler une affirmation sans ambiguïté, d'examiner ses hypothèses et de rédiger une preuve vérifiable. Ce polycopié s'adresse aux étudiants de deuxième année de Licence Mathématiques et accompagne le travail en séance de travaux dirigés.

L'ouvrage suit une progression en quatre chapitres : langage et raisonnement mathématiques ; théorie des ensembles ; calcul propositionnel et calcul des prédicats ; bon ordre et récurrence. Les quatre chapitres sont accompagnés d'exemples guidés, d'illustrations TikZ et d'exercices progressifs allant jusqu'au niveau « Défi ». Les résultats d'indépendance et les théorèmes les plus avancés sont présentés avec leurs hypothèses et la portée exacte de leurs énoncés. Les énoncés et exercices sont rédigés pour être travaillés activement : rechercher des exemples et des contre-exemples, identifier les quantificateurs, puis expliquer chaque étape d'une démonstration. La version augmentée approfondit les graphes d'applications, les paradoxes autoréférentiels, l'antisymétrie et deux méthodes de récurrence illustrées par des inégalités classiques.
\tableofcontents
\mainmatter
\chapter{Langage et raisonnement mathématiques}
\label{chap:langage}
\begin{flushleft}\small\sffamily
\textbf{Prérequis :} calcul algébrique élémentaire, ensembles et fonctions vus en L1.\qquad
\textbf{Volume indicatif :} trois semaines (cours et TD).
\end{flushleft}

\begin{questionfondatrice}
Comment passer d'une observation convaincante à une affirmation précisément formulée, puis à une démonstration dont chaque étape peut être vérifiée ?
\end{questionfondatrice}
\begin{objectifs}
À l'issue du chapitre, l'étudiant doit savoir distinguer énoncé, hypothèse et conclusion ; reconnaître définition, axiome, conjecture et théorème ; analyser les mots \emph{tout}, \emph{il existe} et \emph{si\ldots alors} ; rédiger une preuve directe, par contraposée, par l'absurde ou par disjonction des cas ; construire un contre-exemple ; expliquer les conditions d'emploi de « sans perte de généralité ».
\end{objectifs}

\section{Pourquoi apprendre à démontrer ?}
Une suite de valeurs concordantes ne constitue pas une preuve d'un énoncé universel. Ainsi, les valeurs de $n^2+n+41$ sont premières pour $n=0,1,\ldots,39$, mais $n=40$ donne $40^2+40+41=41^2$, qui n'est pas premier. Inversement, vérifier l'égalité $1+3+\cdots+(2n-1)=n^2$ pour quelques entiers n'explique pas pourquoi elle doit rester vraie pour tout $n$.

Le cours de logique répond à deux questions différentes : \emph{que signifie exactement l'affirmation ?} et \emph{qu'est-ce qui autorise à conclure ?} Une définition fixe le sens des mots ; une démonstration établit un lien contrôlable entre hypothèses et conclusion.

\begin{exemple}{Trois niveaux d'information}{trois-niveaux}
On lit : « Le carré d'un entier impair est impair. »
\begin{enumerate}
\item \textbf{Observation :} $1^2=1$, $3^2=9$ et $5^2=25$ sont impairs.
\item \textbf{Formulation :} pour tout $n\in\Z$, si $n$ est impair, alors $n^2$ est impair.
\item \textbf{Preuve :} il existe $k\in\Z$ tel que $n=2k+1$ ; ainsi $n^2=2(2k^2+2k)+1$.
\end{enumerate}
Ce dernier calcul prouve l'affirmation pour \emph{chaque} entier impair et non pour les seuls exemples examinés.
\end{exemple}

\begin{figure}[htbp]\centering
\begin{tikzpicture}[>=Latex,node distance=8mm,
 box/.style={draw=LFblue,fill=LFpale,rounded corners=2mm,align=center,text width=4cm,minimum height=13mm,font=\small}]
\node[box] (obs) {Observation\\exemples et calculs};
\node[box,right=13mm of obs] (conj) {Conjecture\\énoncé précis};
\node[box,below=12mm of conj] (pr) {Démonstration\\argument général};
\node[box,left=13mm of pr] (th) {Théorème\\résultat établi};
\draw[->,thick,LFteal] (obs)--(conj);
\draw[->,thick,LFteal] (conj)--(pr);
\draw[->,thick,LFteal] (pr)--(th);
\draw[->,thick,LFcoral,dashed] (conj.west) to[bend right=22] node[above,font=\footnotesize]{contre-exemple} (obs.south);
\end{tikzpicture}
\caption{L'exploration suggère un résultat ; seule une preuve le fonde, et un contre-exemple peut le réfuter.}\label{fig:log:demarche}
\end{figure}

\section{Le vocabulaire d'un texte mathématique}
\begin{definition}{Énoncé et proposition}{enonce}
Une \emph{proposition} est un énoncé auquel on peut attribuer une valeur de vérité, vrai ou faux, une fois son domaine et les termes qui y figurent fixés. Une phrase ouverte, telle que « $x^2>4$ », devient une proposition lorsqu'on fixe $x$ ou qu'on quantifie la variable.
\end{definition}
\begin{exemple}{L'importance du domaine}{domaine}
« Tout nombre positif possède une racine carrée » est vrai dans $\R$ avec la convention $x\geq0$ ; la formulation « tout nombre réel possède une racine carrée réelle » est fausse. L'énoncé « $x^2=2$ a une solution » est vrai dans $\R$ et faux dans $\Q$. Le domaine n'est donc jamais un détail.
\end{exemple}

\begin{definition}{Définition, axiome, conjecture, théorème}{statut}
Une \emph{définition} introduit un terme ou une notation ; un \emph{axiome} est un énoncé admis dans un système donné ; une \emph{conjecture} est une affirmation envisagée mais non encore démontrée ; un \emph{théorème} est établi par une démonstration à partir des règles et résultats admis. On appelle souvent \emph{lemme} un résultat auxiliaire et \emph{corollaire} une conséquence proche d'un résultat déjà démontré.
\end{definition}
\begin{attention}
Une définition n'est pas à « démontrer » : on vérifie plutôt qu'un objet satisfait ou non ses conditions. Un théorème, lui, appelle une justification. Un axiome n'est pas une vérité magique : son rôle dépend du système mathématique choisi.
\end{attention}

\subsection{Hypothèses et conclusion}
La forme classique d'un théorème est : « Soit $x$ un objet satisfaisant $H(x)$. Alors $C(x)$. » Pour rédiger sa preuve, commencer par annoncer ce qui est supposé et ce qui doit être obtenu. Dans la phrase « Si $n$ est multiple de $6$, alors $n$ est pair », l'hypothèse est « $n$ multiple de $6$ », la conclusion « $n$ pair ».

\begin{exemple}{La réciproque n'est pas la proposition}{reciproque}
L'implication « multiple de $6$ $\Rightarrow$ pair » est vraie. Sa réciproque « pair $\Rightarrow$ multiple de $6$ » est fausse : $n=2$ est un contre-exemple. Le fait qu'une implication soit vraie ne donne aucune permission de la retourner.
\end{exemple}

\section{Traduire une phrase en langage mathématique}
Le calcul formel des propositions et des prédicats sera développé au chapitre 3. Nous introduisons ici seulement les notations utiles à la lecture et à la rédaction d'une preuve.

\begin{center}\renewcommand{\arraystretch}{1.3}
\begin{tabularx}{\textwidth}{>{\bfseries}l l X}\toprule
Expression & Notation & Sens dans un énoncé \\ \midrule
pour tout & $\forall$ & chaque élément du domaine satisfait la propriété \\
il existe & $\exists$ & au moins un élément la satisfait \\
si\ldots alors & $\Rightarrow$ & l'hypothèse entraîne la conclusion \\
si et seulement si & $\Leftrightarrow$ & deux implications à établir \\
non & $\neg$ & négation de l'énoncé \\
et, ou & $\wedge,\ \vee$ & « ou » inclusif, sauf indication contraire \\ \bottomrule
\end{tabularx}\end{center}

\begin{exemple}{Traduire sans ambiguïté}{traduire}
« Toute fonction continue en $a$ est bornée au voisinage de $a$ » peut se lire : pour toute fonction $f$ et tout point $a$ de son domaine où elle est continue, il existe $\delta>0$ et $M>0$ tels que $|f(x)|\le M$ pour tout $x$ du domaine vérifiant $|x-a|<\delta$. Les mots « il existe » viennent \emph{après} la fonction et le point fixés : $M$ et $\delta$ peuvent en dépendre.
\end{exemple}

\subsection{Le rôle de l'ordre des quantificateurs}
Comparer les deux phrases :
\[
\forall x\in\R,\ \exists y\in\R : y>x,
\qquad
\exists y\in\R,\ \forall x\in\R : y>x.
\]
La première est vraie (choisir $y=x+1$), la seconde est fausse (choisir $x=y+1$). L'ordre des quantificateurs modifie le sens.

\begin{figure}[htbp]\centering
\begin{tikzpicture}[>=Latex]
\draw[->,LFnavy,thick] (-.1,0)--(10.4,0) node[right] {$\R$};
\foreach \x in {1.5,4,6.5,9} {\draw (\x,.10)--(\x,-.10);}
\fill[LFblue] (4,0) circle(2pt) node[below=4pt] {$x$};
\fill[LFteal] (6.5,0) circle(2pt) node[below=4pt] {$y=x+1$};
\draw[->,LFteal,thick] (4,.35)--(6.5,.35) node[midway,above,font=\small] {choix dépendant de $x$};
\end{tikzpicture}
\caption{« Pour chaque $x$, il existe $y$ » autorise un choix de $y$ différent pour chaque $x$.}\label{fig:log:quantif}
\end{figure}

\begin{proposition}{Négation des énoncés quantifiés}{negation}
Pour une propriété $P(x)$ sur un domaine $E$, on a
\[
\neg(\forall x\in E,\,P(x))\ \Leftrightarrow\ \exists x\in E,\,\neg P(x),
\]
\[
\neg(\exists x\in E,\,P(x))\ \Leftrightarrow\ \forall x\in E,\,\neg P(x).
\]
\end{proposition}
\begin{proof}
Dire que $P$ ne vaut pas pour tous les éléments signifie qu'il en existe au moins un pour lequel elle échoue. Dire qu'il n'existe aucun élément satisfaisant $P$ signifie qu'aucun n'y satisfait. Ce sont les définitions du « pour tout » et du « il existe ».
\end{proof}
\begin{exemple}{Négation d'une limite}{neg-lim}
La négation de « pour tout $\varepsilon>0$, il existe $N$ tel que pour tout $n\ge N$, $|u_n-\ell|<\varepsilon$ » est :
\[
\exists\varepsilon_0>0,\ \forall N\in\N,\ \exists n\ge N:\ |u_n-\ell|\ge\varepsilon_0.
\]
On renverse chaque quantificateur et on nie seulement la propriété finale.
\end{exemple}

\section{Les méthodes classiques de démonstration}
\subsection{La preuve directe}
On suppose l'hypothèse et l'on en tire la conclusion par une chaîne d'arguments justifiés.
\begin{proposition}{Somme de deux entiers pairs}{pairs}
Si $a,b\in\Z$ sont pairs, alors $a+b$ est pair.
\end{proposition}
\begin{proof}
Il existe $k,\ell\in\Z$ tels que $a=2k$ et $b=2\ell$. Alors $a+b=2(k+\ell)$ avec $k+\ell\in\Z$ : la somme est paire par définition.
\end{proof}
\begin{retenir}
Dans une preuve directe, chaque ligne doit venir d'une hypothèse, d'une définition ou d'un résultat déjà disponible. Les calculs n'ont de valeur probante que si leur rôle est explicité.
\end{retenir}

\subsection{La contraposée}
Pour établir $P\Rightarrow Q$, on peut établir l'implication équivalente $\neg Q\Rightarrow\neg P$.
\begin{proposition}{Le carré pair impose la parité}{carre-pair}
Si $n\in\Z$ et $n^2$ est pair, alors $n$ est pair.
\end{proposition}
\begin{proof}
Supposons que $n$ ne soit pas pair : il existe $k\in\Z$ tel que $n=2k+1$. Dès lors $n^2=2(2k^2+2k)+1$ est impair. C'est la contraposée cherchée.
\end{proof}
\begin{attention}
La \emph{contraposée} de $P\Rightarrow Q$ est $\neg Q\Rightarrow\neg P$, pas la \emph{réciproque} $Q\Rightarrow P$. La première est logiquement équivalente à l'énoncé initial ; la seconde ne l'est pas en général.
\end{attention}

\subsection{Le raisonnement par l'absurde}
Pour démontrer $P$, on suppose $\neg P$ et l'on en déduit une contradiction. Le raisonnement par l'absurde est particulièrement utile pour l'impossibilité et l'unicité.
\begin{theoreme}{Irrationalité de $\sqrt2$}{sqrt2}
Le nombre $\sqrt2$ n'appartient pas à $\Q$.
\end{theoreme}
\begin{proof}
Supposons $\sqrt2=p/q$ avec $p,q\in\Z$, $q\ne0$ et $p,q$ premiers entre eux. Alors $p^2=2q^2$ ; par la proposition précédente, $p$ est pair : $p=2r$. Ainsi $q^2=2r^2$, donc $q$ est pair. Les nombres $p$ et $q$ sont tous deux divisibles par $2$, contradiction avec l'irréductibilité de la fraction.
\end{proof}

\begin{figure}[htbp]\centering
\begin{tikzpicture}[>=Latex,node distance=8mm,
 b/.style={draw=LFblue,rounded corners=2mm,fill=LFpale,minimum width=42mm,align=center,minimum height=11mm,font=\small}]
\node[b] (a) {Supposer $\sqrt2=p/q$\\ fraction irréductible};
\node[b,below=of a] (b) {$p^2=2q^2$\\donc $p$ pair};
\node[b,below=of b] (c) {$q^2=2r^2$\\donc $q$ pair};
\node[b,below=of c,draw=LFcoral,fill=LFred] (d) {$2$ divise $p$ et $q$\\contradiction};
\foreach \x/\y in {a/b,b/c,c/d}{\draw[->,LFteal,thick] (\x)--(\y);}
\end{tikzpicture}
\caption{Architecture d'une preuve par l'absurde : hypothèse niée, conséquences, contradiction.}\label{fig:log:absurde}
\end{figure}

\subsection{La disjonction des cas}
Lorsque les hypothèses couvrent plusieurs situations distinctes, il suffit de conclure dans chacun des cas.
\begin{proposition}{Le produit de deux entiers consécutifs}{consecutifs}
Pour tout $n\in\Z$, le produit $n(n+1)$ est pair.
\end{proposition}
\begin{proof}
Si $n$ est pair, le produit est pair. Si $n$ est impair, $n+1$ est pair et le produit est encore pair. Ces deux cas couvrent tous les entiers.
\end{proof}

\subsection{Démontrer une équivalence}
Établir $P\Leftrightarrow Q$ exige en général les deux sens : $P\Rightarrow Q$ puis $Q\Rightarrow P$.
\begin{proposition}{Caractérisation d'une inclusion}{inclu}
Pour deux parties $A,B$ d'un ensemble $E$, on a $A\subseteq B$ si et seulement si $A\cap B=A$.
\end{proposition}
\begin{proof}
Si $A\subseteq B$, tout élément de $A$ appartient à $B$, donc $A\cap B=A$. Réciproquement, si $A\cap B=A$ et $x\in A$, alors $x\in A\cap B$, donc $x\in B$ ; ainsi $A\subseteq B$.
\end{proof}

\section{Existence, construction et unicité}
Une preuve d'existence doit établir qu'au moins un objet satisfait les conditions. Elle peut être \emph{constructive} (on présente l'objet) ou non constructive (on prouve l'existence sans le déterminer explicitement). Une preuve d'unicité part de deux objets satisfaisant la propriété et montre qu'ils sont égaux.

\begin{exemple}{Existence constructive}{existence}
Pour chaque $a\in\R$, l'équation $x+3=a$ possède une solution : choisir $x=a-3$ et vérifier $x+3=a$. Si $x$ et $y$ sont solutions, alors $x+3=y+3$, d'où $x=y$ : la solution est unique.
\end{exemple}
\begin{proposition}{Unicité d'une limite réelle}{unique-lim}
Une suite réelle ne peut pas avoir deux limites distinctes.
\end{proposition}
\begin{proof}
Supposons $u_n\to\ell$ et $u_n\to m$. Si $\ell\ne m$, posons $\varepsilon=|\ell-m|/3>0$. Pour $n$ assez grand, $|u_n-\ell|<\varepsilon$ et $|u_n-m|<\varepsilon$, donc $|\ell-m|\le|\ell-u_n|+|u_n-m|<2|\ell-m|/3$, contradiction. Ainsi $\ell=m$.
\end{proof}

\section{Contre-exemples et vigilance sur la généralité}
Un énoncé universel est réfuté par un seul objet qui satisfait les hypothèses mais pas la conclusion. Il ne suffit pas de citer un objet hors du domaine ou qui viole une hypothèse.
\begin{contreexemple}{Une implication fausse}{faux}
« Si $ab=0$, alors $a=b=0$ » est faux sur $\R$ : prendre $a=0$, $b=1$. Le contre-exemple satisfait bien $ab=0$, mais ne satisfait pas $a=b=0$. La conclusion correcte est $a=0$ \emph{ou} $b=0$.
\end{contreexemple}
\begin{contreexemple}{La convergence n'autorise pas tout}{faux-lim}
« Si $u_n\to0$, alors $1/u_n\to0$ » est faux : $u_n=1/n$ pour $n\ge1$ converge vers $0$, mais $1/u_n=n$ diverge vers $+\infty$. La valeur $u_n=0$ doit, de plus, être exclue pour définir l'inverse.
\end{contreexemple}

\subsection{Quand peut-on dire « sans perte de généralité » ?}
Cette formule est correcte seulement si un argument de symétrie ou de réduction permet de passer du cas général au cas choisi.
\begin{exemple}{Une symétrie justifiée}{sptg}
Pour $a,b\in\R$, l'égalité $\max(a,b)+\min(a,b)=a+b$ se démontre en supposant, sans perte de généralité, $a\le b$ : les rôles de $a$ et $b$ peuvent être échangés, et les deux membres sont symétriques. Dans ce cas, $\min(a,b)=a$ et $\max(a,b)=b$.
\end{exemple}
\begin{attention}
On ne peut pas, sans justification, supposer « $x>0$ » dans une proposition portant sur tous les réels. Pour la fonction $f(x)=x^3$, les comportements en $x>0$ et $x<0$ ne sont pas automatiquement interchangeables sans exploiter une propriété précise, par exemple l'imparité.
\end{attention}

\section{Comment rédiger une preuve lisible ?}
Une démonstration doit nommer les objets, citer les hypothèses pertinentes, annoncer l'objectif et motiver les transitions. Elle doit éviter les enchaînements de symboles qui ne précisent pas si l'on effectue une implication, une équivalence ou une simple transformation.
\begin{retenir}
\textbf{Plan de rédaction conseillé :} (1) fixer le cadre et les objets ; (2) annoncer la méthode si elle n'est pas évidente ; (3) utiliser les définitions ; (4) enchaîner les justifications ; (5) conclure exactement à l'énoncé demandé. Les mots \emph{donc}, \emph{car}, \emph{en effet} et \emph{réciproquement} marquent la structure de l'argument.
\end{retenir}
\begin{exemple}{Réparer une justification incomplète}{reparer}
Affirmation : si $x^2<4$, alors $-2<x<2$. Écrire seulement « on prend la racine » est ambigu. On écrit plutôt : $x^2<4$ équivaut à $|x|<2$, donc $-2<x<2$ par définition de la valeur absolue.
\end{exemple}

\section{Exercices progressifs}
Les exercices ci-dessous peuvent être utilisés en TD ; ils sont à chercher d'abord seul, puis à discuter en séance.
\subsection*{Niveau 1 — Lire et formuler}
\begin{exercice}[Hypothèse et conclusion]\label{ex:log:1}
Identifier l'hypothèse et la conclusion : « Si $x\in\R$ et $x>2$, alors $x^2>4$. » Écrire la réciproque et dire si elle est vraie.
\end{exercice}
\begin{exercice}[Négations]\label{ex:log:2}
Nier exactement : (a) $\forall x\in\R,\ x^2\ge0$ ; (b) $\exists n\in\N,\ n^2=2$ ; (c) $\forall\varepsilon>0\ \exists N\ \forall n\ge N,\ |u_n|<\varepsilon$.
\end{exercice}
\begin{exercice}[Quantificateurs]\label{ex:log:3}
Déterminer la valeur de vérité de $\forall x\in\R\ \exists y\in\R:\ x+y=0$ et de $\exists y\in\R\ \forall x\in\R:\ x+y=0$.
\end{exercice}
\begin{exercice}[Définition et théorème]\label{ex:log:4}
Classer les phrases suivantes : « un entier est pair s'il s'écrit $2k$ » ; « la somme de deux entiers pairs est paire » ; « pour tout entier $n$, $n^2+n+41$ est premier ».
\end{exercice}
\subsection*{Niveau 2 — Prouver et réfuter}
\begin{exercice}[Preuve directe]\label{ex:log:5}
Montrer que la somme de deux entiers impairs est paire.
\end{exercice}
\begin{exercice}[Contraposée]\label{ex:log:6}
Montrer : si $n^2$ est impair, alors $n$ est impair.
\end{exercice}
\begin{exercice}[Disjonction des cas]\label{ex:log:7}
Montrer que $n^2+n$ est pair pour tout $n\in\Z$.
\end{exercice}
\begin{exercice}[Équivalence]\label{ex:log:8}
Pour $x\in\R$, montrer que $|x|<1$ si et seulement si $-1<x<1$.
\end{exercice}
\begin{exercice}[Contre-exemples]\label{ex:log:9}
Pour chacune des assertions suivantes, donner un contre-exemple ou une preuve : (a) $a^2=b^2\Rightarrow a=b$ ; (b) $a\le b\Rightarrow a^2\le b^2$ ; (c) $a,b>0\Rightarrow ab>0$.
\end{exercice}
\subsection*{Niveau 3 — Construire un argument}
\begin{exercice}[Existence et unicité]\label{ex:log:10}
Montrer que, pour tout $a\in\R$, l'équation $3x-4=a$ a une unique solution réelle.
\end{exercice}
\begin{exercice}[L'absurde]\label{ex:log:11}
Démontrer qu'il n'existe pas de plus grand entier naturel.
\end{exercice}
\begin{exercice}[Deux sens]\label{ex:log:12}
Pour $A,B\subseteq E$, montrer que $A\subseteq B$ si et seulement si $A\cup B=B$.
\end{exercice}
\begin{exercice}[Un « sans perte de généralité »]\label{ex:log:13}
Justifier précisément la possibilité de supposer $a\le b$ pour démontrer $|a-b|=\max(a,b)-\min(a,b)$ ; terminer la preuve.
\end{exercice}
\begin{exercice}[Réfuter une preuve]\label{ex:log:14}
Un étudiant écrit : « $n=1,2,3,4$ vérifie $n^2+n+41$ premier, donc c'est vrai pour tout $n\in\N$. » Expliquer la faute logique et donner un contre-exemple explicite.
\end{exercice}
\begin{exercice}[Énoncé quantifié]\label{ex:log:15}
Écrire formellement puis nier : « Toute fonction $f:\R\to\R$ admet au moins un zéro. » Dire si la phrase initiale est vraie.
\end{exercice}

\subsection*{Niveau 5 — Défis}
\begin{exercice}[Aucun entier entre deux consécutifs]\label{ex:log:16}
Montrer qu'il n'existe pas d'entier $k$ tel que $n<k<n+1$ pour un $n\in\Z$ fixé. Identifier l'hypothèse mathématique utilisée.
\end{exercice}
\begin{exercice}[Défi des quantificateurs]\label{ex:log:17}
Écrire formellement et nier : « Pour toute suite réelle bornée, il existe un majorant réel ». Indiquer pourquoi la négation est fausse.
\end{exercice}
\begin{exercice}[Une preuve à réparer]\label{ex:log:18}
On prétend : « $ab=0$ implique $a=0$ puisque l'on divise par $b$. » Expliquer l'erreur et donner une démonstration correcte dans $\R$ de $ab=0\Rightarrow(a=0\text{ ou }b=0)$.
\end{exercice}

\begin{retenir}
Le chapitre suivant étudiera de manière systématique les ensembles, les relations et les applications ; le calcul propositionnel et celui des prédicats formaliseront ensuite les structures de raisonnement rencontrées ici.
\end{retenir}

\chapter{Théorie des ensembles}
\label{chap:ensembles}
\begin{flushleft}\small\sffamily
\textbf{Prérequis :} quantificateurs, applications, méthodes de preuve du chapitre 1.\qquad
\textbf{Volume indicatif :} quatre semaines (cours et TD).
\end{flushleft}
\begin{questionfondatrice}
Comment comparer la taille de deux ensembles, même infinis, et pourquoi la formule « l'ensemble de tous les ensembles » ne peut-elle être prise au pied de la lettre ?
\end{questionfondatrice}
\begin{objectifs}
Manipuler les opérations ensemblistes et leurs preuves par appartenance ; distinguer relation et application ; établir la dénombrabilité de $\Z$ et $\Q$ ; comprendre l'argument diagonal et le théorème de Cantor--Bernstein ; identifier exactement le rôle des axiomes dans les résultats d'indépendance.
\end{objectifs}

\section{Ensembles, appartenance et opérations}
Un ensemble $A$ est déterminé par ses éléments. Écrire $x\in A$ signifie que $x$ appartient à $A$ ; écrire $A\subseteq B$ signifie $\forall x\,(x\in A\Rightarrow x\in B)$. Deux ensembles sont égaux exactement quand ils ont les mêmes éléments (\emph{extensionalité}).
\begin{definition}{Opérations fondamentales}{operations}
Pour $A,B\subseteq E$, on pose
\[
A\cup B=\{x:x\in A\ \text{ou}\ x\in B\},\quad
A\cap B=\{x:x\in A\ \text{et}\ x\in B\},\quad
A\setminus B=\{x\in A:x\notin B\}.
\]
Le complémentaire est $A^c=E\setminus A$ (par rapport à $E$). L'ensemble des parties est $\PP(E)=\{A:A\subseteq E\}$.
\end{definition}
\begin{exemple}{Appartenance, inclusion et ensemble des parties}{parties}
Si $E=\{a,b\}$, alors $\PP(E)=\{\varnothing,\{a\},\{b\},E\}$ ; $a\in E$, mais $\{a\}\subseteq E$. Il ne faut pas confondre $a$ et $\{a\}$. Pour $|E|=n$, on a $|\PP(E)|=2^n$ : chaque élément est choisi ou non.
\end{exemple}
\begin{proposition}{Lois de De Morgan}{demorgan}
Pour tous $A,B\subseteq E$, $(A\cup B)^c=A^c\cap B^c$ et $(A\cap B)^c=A^c\cup B^c$. Pour une famille $(A_i)_{i\in I}$, ces égalités valent en remplaçant « deux » par « tous les indices ».
\end{proposition}
\begin{proof}Fixons $x\in E$. On a $x\in(A\cup B)^c$ si et seulement si $x\notin A$ et $x\notin B$, ce qui équivaut à $x\in A^c\cap B^c$. La seconde formule s'obtient de la même manière. L'argument repose sur la négation de « il existe » et de « pour tout » et vaut pour une famille quelconque.\end{proof}
\begin{figure}[htbp]\centering
\begin{tikzpicture}[scale=.87]
\foreach \dx/\lab in {0/{$A\cup B$},5.4/{$A\cap B$}}{
 \begin{scope}[xshift=\dx cm]
 \draw[LFnavy] (-.3,-.7) rectangle (4.6,2.7);
 \fill[LFblue!18] (1.7,1) circle (1.13);
 \fill[LFblue!18] (2.7,1) circle (1.13);
 \draw[LFblue,thick] (1.7,1) circle (1.13) node[above left=7mm] {$A$};
 \draw[LFteal,thick] (2.7,1) circle (1.13) node[above right=7mm] {$B$};
 \node[below] at (2.1,-.85) {\lab};
 \end{scope}}
\begin{scope}[xshift=5.4cm]
\begin{scope}
\clip (1.7,1) circle (1.13);
\fill[LFteal!55] (2.7,1) circle (1.13);
\end{scope}\end{scope}
\end{tikzpicture}
\caption{Une région est décrite par une condition logique sur l'appartenance. La superposition correspond à « et ».}
\end{figure}

\section{Produit cartésien, relations et applications}
\begin{definition}{Produit et relation}{relation}
$A\times B=\{(a,b):a\in A,\ b\in B\}$. Une relation binaire sur $E$ est une partie $R\subseteq E\times E$. Elle est réflexive si $xRx$ pour tout $x$, symétrique si $xRy\Rightarrow yRx$, transitive si $xRy$ et $yRz\Rightarrow xRz$.
\end{definition}
Une \emph{relation d'équivalence} possède ces trois propriétés ; ses classes partitionnent $E$. Une relation d'ordre est réflexive, antisymétrique et transitive. Une application $f:A\to B$ associe à chaque $a\in A$ une unique image $f(a)\in B$ ; l'image de $S\subseteq A$ et l'image réciproque de $T\subseteq B$ sont
\[f(S)=\{f(x):x\in S\},\qquad f^{-1}(T)=\{x\in A:f(x)\in T\}.\]
\begin{attention}
La notation $f^{-1}(T)$ existe pour \emph{toute} application : elle n'exige pas de fonction réciproque $f^{-1}:B\to A$. Les images réciproques commutent aux unions, intersections et compléments ; les images directes commutent aux unions, mais pas nécessairement aux intersections.
\end{attention}
\begin{exemple}{Une intersection perdue par l'image directe}{intersection-image}
Pour $f:\R\to\R$, $f(x)=x^2$, prendre $S=\{-1\}$ et $T=\{1\}$. Alors $S\cap T=\varnothing$ mais $f(S)\cap f(T)=\{1\}$. L'injectivité permet de retrouver l'égalité pour les intersections.
\end{exemple}
\begin{proposition}{Images réciproques}{preimages}
Pour toute famille $(B_i)_{i\in I}$ de parties de $B$,
\[f^{-1}\!\left(\bigcup_i B_i\right)=\bigcup_i f^{-1}(B_i),\qquad
f^{-1}\!\left(\bigcap_i B_i\right)=\bigcap_i f^{-1}(B_i).\]
\end{proposition}
\begin{proof}Un point $x$ appartient à la première partie si et seulement s'il existe $i$ tel que $f(x)\in B_i$. Pour la seconde, remplacer « il existe » par « pour tout ».\end{proof}


\subsection{Une application est une relation particulière}
Une relation de $A$ vers $B$ est un sous-ensemble de $A\times B$. Une application
$f:A\to B$ peut donc être identifiée à son \emph{graphe}
\[
 \Gamma_f=\{(a,b)\in A\times B:b=f(a)\}.
\]
La différence avec une relation quelconque tient à deux exigences, pour chaque $a\in A$ :
\emph{existence} d'une image et \emph{unicité} de cette image.
\begin{proposition}{Caractérisation ensembliste d'une application}{graphe-fonction}
Une partie $G\subseteq A\times B$ est le graphe d'une application de $A$ vers $B$ si et seulement si
\[
 \forall a\in A\;\exists!b\in B:\ (a,b)\in G.
\]
Pour un $a$ fixé, sa fibre verticale $G\cap(\{a\}\times B)$ contient donc exactement un point.
\end{proposition}
\begin{proof}
Si $G=\Gamma_f$, le couple $(a,f(a))$ appartient à $G$ et tout autre couple $(a,b)$ de $G$ impose $b=f(a)$.
Réciproquement, la propriété définit sans ambiguïté $f(a)$ comme l'unique $b$ tel que $(a,b)\in G$. On retrouve alors $G=\Gamma_f$.
\end{proof}
\begin{exemple}{Un graphe qui n'est pas celui d'une fonction}{graphe-exemple}
Sur $A=\{1,2\}$ et $B=\{u,v\}$, la relation
$G_1=\{(1,u),(2,v)\}$ définit une application. La relation
$G_2=\{(1,u),(1,v),(2,u)\}$ ne le fait pas : $1$ a deux images.
La relation $G_3=\{(1,u)\}$ ne le fait pas davantage : $2$ n'a aucune image.
\end{exemple}
\begin{figure}[htbp]\centering
\begin{tikzpicture}[>=Latex,scale=.92]
\foreach \d/\key/\head in {0/L/{graphe d'une application},5.4/R/{deux images pour $1$}}{
 \node[font=\small\sffamily\bfseries,LFnavy] at (\d+1.7,2.15) {\head};
 \draw[LFblue,thick] (\d,0) ellipse (.53 and 1.15);
 \draw[LFblue,thick] (\d+3.2,0) ellipse (.53 and 1.15);
 \node at (\d,.53) (a\key) {$1$};\node at (\d,-.55) (b\key) {$2$};
 \node at (\d+3.2,.53) (u\key) {$u$};\node at (\d+3.2,-.55) (v\key) {$v$};
}
\draw[->,LFteal,thick] (aL)--(uL); \draw[->,LFteal,thick] (bL)--(vL);
\draw[->,LFcoral,thick] (aR)--(uR);\draw[->,LFcoral,thick] (aR)--(vR);
\draw[->,LFteal,thick] (bR)--(uR);
\end{tikzpicture}
\caption{Le critère graphique : une et une seule flèche doit partir de chaque élément du domaine.}
\end{figure}


\section{Le paradoxe de Russell et les règles de formation}
La théorie naïve permet de décrire des collections par propriétés. Cette liberté absolue aboutit à une contradiction. Supposons qu'il existe un ensemble $R=\{x:x\notin x\}$. Alors
\[R\in R\quad\Longleftrightarrow\quad R\notin R.\]
Les deux possibilités sont contradictoires. Le problème est le postulat implicite que \emph{toute} propriété définirait un ensemble. En théorie axiomatique ZF, le schéma de séparation ne fabrique que $\{x\in E:P(x)\}$ à partir d'un ensemble $E$ déjà donné ; il ne postule pas un ensemble universel.
\begin{retenir}
Le paradoxe ne prouve pas que les mathématiques sont incohérentes : il montre qu'une règle naïve de formation des ensembles est trop générale. Les axiomes précisent quelles constructions sont admises.
\end{retenir}


\subsection{Le bibliothécaire et le menteur : distinguer les mécanismes}
\begin{exemple}{Le bibliothécaire : une incarnation de Russell}{bibliothecaire}
Dans une bibliothèque fictive, un bibliothécaire devrait répertorier exactement tous les
catalogues qui ne se répertorient pas eux-mêmes. Si son propre catalogue $C$ devait être
catalogué selon cette règle, alors $C$ se répertorie lui-même si et seulement s'il ne se
répertorie pas lui-même. La contradiction reproduit la forme $R\in R\Leftrightarrow R\notin R$.
La difficulté ne vient pas d'une bibliothèque réelle : c'est la prescription
\emph{sans restriction} qui est impossible à satisfaire.
\end{exemple}
\begin{exemple}{Le menteur : une autre famille de paradoxe}{menteur}
La phrase « Cette phrase est fausse » est autoréférentielle. Si l'on admet qu'elle possède
une valeur de vérité classique, cette valeur se renverse aussitôt. Elle ressemble à Russell
par l'autoréférence, mais n'est \emph{pas} une définition ensembliste ; elle appartient aux
paradoxes sémantiques. La version « tous les Crétois sont menteurs », attribuée à un Crétois,
ne produit pas, à elle seule, une contradiction formelle : il faut préciser le sens de « menteur »
et les hypothèses. On ne confondra donc pas paradoxe ensembliste et phrase ambiguë.
\end{exemple}
\begin{attention}
Le paradoxe de Russell concerne la formation illimitée d'ensembles ; le menteur concerne la
vérité d'une phrase parlant de sa propre vérité. Une analogie de structure ne permet pas de
substituer l'un à l'autre dans une démonstration.
\end{attention}


\section{Équipotence et ensembles dénombrables}
\begin{definition}{Équipotence et dénombrabilité}{equipotence}
Deux ensembles $A,B$ sont équipotents, noté $A\sim B$, s'il existe une bijection $A\to B$. Un ensemble est \emph{dénombrable} s'il est fini ou s'il est équipotent à $\N$ (on dit aussi \emph{au plus dénombrable}). Il est \emph{dénombrablement infini} s'il est équipotent à $\N$.
\end{definition}
\begin{exemple}{Une partie propre peut avoir la même taille}{pairs-infinis}
La fonction $n\mapsto2n$ donne une bijection de $\N$ sur les entiers pairs naturels. Pour l'infini, une inclusion stricte n'implique donc pas une cardinalité strictement plus petite.
\end{exemple}
\begin{proposition}{Les entiers relatifs et les rationnels sont dénombrables}{q-denom}
Les ensembles $\Z$ et $\Q$ sont dénombrables.
\end{proposition}
\begin{proof}
Énumérer $\Z$ selon $0,1,-1,2,-2,\ldots$. Pour $\Q$, chaque rationnel est $p/q$ avec $p\in\Z$ et $q\in\N^*$. Les couples $(p,q)$ s'énumèrent par diagonales $|p|+q=1,2,\ldots$, en ignorant les doublons des écritures fractionnaires. Cette énumération surjective montre que $\Q$ est au plus dénombrable ; comme il est infini, il est équipotent à $\N$.
\end{proof}
\begin{figure}[htbp]\centering
\begin{tikzpicture}[x=1.0cm,y=.72cm,>=Latex]
\foreach \i in {1,...,5}{\foreach \j in {1,...,5}{\fill[LFblue] (\i,-\j) circle(1.5pt);}}
\foreach \i in {1,...,5}{\node[above,font=\small] at (\i,-.7) {$\i$};\node[left,font=\small] at (.73,-\i) {$\i$};}
\draw[->,LFteal,very thick,rounded corners=3pt] (1,-1)--(1,-2)--(2,-1)--(3,-1)--(2,-2)--(1,-3)--(1,-4)--(2,-3)--(3,-2)--(4,-1);
\node[anchor=west,font=\small] at (5.6,-2.1) {énumération par};\node[anchor=west,font=\small] at (5.6,-2.55) {diagonales finies};
\end{tikzpicture}
\caption{Chaque couple d'entiers positifs finit par apparaître sur une diagonale.}
\end{figure}
\begin{proposition}{Union dénombrable}{union-denom}
Une union dénombrable d'ensembles au plus dénombrables est au plus dénombrable.
\end{proposition}
\begin{proof}
Pour chaque ensemble non vide $A_n$, choisir une énumération (avec répétitions si nécessaire) $a_{n,k}$. Énumérer les couples $(n,k)$ par diagonales, puis effacer les doublons. Pour une famille dénombrable arbitraire, le choix simultané de ces énumérations utilise une forme dénombrable de l'axiome du choix ; dans les applications concrètes comme $\Q$, les énumérations sont explicites.
\end{proof}

\section{Cantor : tous les infinis n'ont pas la même taille}
\begin{theoreme}{Théorème de Cantor}{cantor}
Pour tout ensemble $E$, il n'existe aucune surjection $f:E\to\PP(E)$. En particulier $|E|<|\PP(E)|$ au sens des cardinalités.
\end{theoreme}
\begin{proof}
Supposons $f$ surjective et formons $D=\{x\in E:x\notin f(x)\}\subseteq E$. Il existerait $d\in E$ tel que $f(d)=D$. Or
\[d\in D\ \Longleftrightarrow\ d\notin f(d)=D,\]
contradiction. Il existe en revanche toujours une injection $E\to\PP(E)$ donnée par $x\mapsto\{x\}$.
\end{proof}
\begin{theoreme}{Diagonalisation sur $[0,1]$}{diagonale}
L'intervalle $[0,1]$ n'est pas dénombrable.
\end{theoreme}
\begin{proof}
Supposons qu'on puisse l'énumérer $x_1,x_2,\ldots$. Écrivons chaque $x_n$ sous sa représentation décimale qui ne se termine pas par une queue infinie de $9$. Construisons $y=0,b_1b_2\ldots$ en posant $b_n=1$ si le $n$-ième chiffre de $x_n$ n'est pas $1$, et $b_n=2$ sinon. Ainsi $b_n\in\{1,2\}$ et diffère du chiffre diagonal de $x_n$. Notre écriture de $y$ ne finit pas par des $9$ et $y\ne x_n$ pour tout $n$ ; contradiction.
\end{proof}
\begin{attention}
La convention sur les décimales est indispensable : $0,5000\ldots=0,4999\ldots$. En choisissant pour $y$ uniquement les chiffres $1$ et $2$, l'ambiguïté disparaît.
\end{attention}
\begin{exemple}{La densité n'implique pas la même cardinalité}{dense-cardinal}
$\Q$ est dense dans $\R$ (tout intervalle ouvert contient un rationnel) mais $\Q$ est dénombrable et $\R$ ne l'est pas. Densité topologique et cardinalité mesurent deux propriétés distinctes.
\end{exemple}

\section{Théorème de Cantor--Bernstein}
\begin{theoreme}{Cantor--Bernstein--Schröder}{cbs}
S'il existe une injection $f:A\to B$ et une injection $g:B\to A$, alors il existe une bijection de $A$ sur $B$.
\end{theoreme}
\begin{proof}
Définissons $A_0=A\setminus g(B)$, puis $A_{n+1}=g(f(A_n))$ et $A_*=\bigcup_{n\ge0}A_n$. On a $g(f(A_*))=A_*\setminus A_0$. Posons
\[
h(a)=\begin{cases}f(a),&a\in A_*,\\g^{-1}(a),&a\in A\setminus A_* .\end{cases}
\]
La seconde branche est définie car $A\setminus A_*\subseteq g(B)$ et $g$ est injective. Les deux images sont disjointes : si $f(a)=g^{-1}(a')$, avec $a\in A_*$ et $a'\notin A_*$, alors $a'=g(f(a))\in A_*$, contradiction. Les branches sont injectives. Enfin, si $b\notin f(A_*)$, alors $g(b)\notin A_*$ (sinon $g(b)\in A_*\setminus A_0=g(f(A_*))$ et $b\in f(A_*)$), donc $h(g(b))=b$. Ainsi $h$ est bijective.
\end{proof}
\begin{exemple}{Une application aux intervalles}{intervalle-card}
L'injection $(0,1)\hookrightarrow[0,1]$ est l'inclusion, et $x\mapsto\frac14+\frac12x$ injecte $[0,1]$ dans $(0,1)$. Cantor--Bernstein donne une bijection sans devoir en écrire laborieusement une formule.
\end{exemple}

\section{Axiome du choix et frontières de la théorie}
\begin{definition}{Axiome du choix (AC)}{choix}
Pour toute famille $(E_i)_{i\in I}$ d'ensembles non vides, il existe une fonction $c$ telle que $c(i)\in E_i$ pour chaque $i$. Le choix finitiste ne pose pas de difficulté ; AC affirme l'existence simultanée pour une famille arbitraire.
\end{definition}
\begin{proposition}{Trois formulations équivalentes dans ZF}{choix-equivalence}
L'axiome du choix, le théorème du bon ordre (tout ensemble admet un bon ordre) et le lemme de Zorn sont équivalents dans la théorie ZF. Le résultat et ses applications seront repris au chapitre 4.
\end{proposition}
L'\emph{hypothèse du continu} (CH) affirme qu'il n'existe pas de cardinal strictement compris entre celui de $\N$ et celui de $\R$ : $2^{\aleph_0}=\aleph_1$. Ce n'est pas un théorème de ZFC. Plus précisément, sous l'hypothèse de cohérence de ZFC, Gödel a construit un modèle où CH et AC sont vrais (univers constructible), et Cohen un modèle de ZFC où CH est faux (méthode du forcing). Ainsi ZFC ne décide pas CH, si ZFC est cohérente. Nous donnons ici la portée de ces résultats, non leur preuve technique, qui exige une formation spécialisée.
\begin{retenir}
Une assertion « indépendante de ZFC » signifie que ni elle ni sa négation ne sont démontrables à partir de ces axiomes, \emph{sous la condition de cohérence appropriée}. Il ne s'agit pas de dire qu'elle est simultanément vraie et fausse dans un même modèle.
\end{retenir}

\section{Exercices progressifs}
\subsection*{Niveau 1 — Manipuler}
\begin{exercice}[Ensemble des parties]\label{ex:ens:1}Écrire $\PP(\{1,2,3\})$ et vérifier son cardinal.\end{exercice}
\begin{exercice}[Égalité par double inclusion]\label{ex:ens:2}Montrer $A\setminus(B\cup C)=(A\setminus B)\cap(A\setminus C)$.\end{exercice}
\begin{exercice}[Images réciproques]\label{ex:ens:3}Pour $f(x)=x^2$, calculer $f^{-1}([1,4])$, $f^{-1}((-\infty,0))$ et $f([-2,1])$.\end{exercice}
\begin{exercice}[Équivalence]\label{ex:ens:4}Sur $\Z$, poser $a\sim b$ si $3\mid(a-b)$. Vérifier que c'est une équivalence et décrire ses classes.\end{exercice}
\subsection*{Niveau 2 — Démontrer}
\begin{exercice}[Distributivité]\label{ex:ens:5}Montrer $A\cap(B\cup C)=(A\cap B)\cup(A\cap C)$ sans diagramme.\end{exercice}
\begin{exercice}[Image et injectivité]\label{ex:ens:6}Prouver que $f(S\cap T)\subseteq f(S)\cap f(T)$ et que l'égalité vaut si $f$ est injective.\end{exercice}
\begin{exercice}[Infinis explicites]\label{ex:ens:7}Construire une bijection de $\N$ vers $\Z$ et une de $\N$ vers $\N\times\N$.\end{exercice}
\begin{exercice}[Sous-ensemble d'un dénombrable]\label{ex:ens:8}Montrer qu'une partie infinie de $\N$ est équipotente à $\N$.\end{exercice}
\subsection*{Niveau 3 — Structurer}
\begin{exercice}[Applications et partitions]\label{ex:ens:9}Montrer que les classes d'une relation d'équivalence constituent une partition de $E$ et réciproquement.\end{exercice}
\begin{exercice}[Cardinalité des parties]\label{ex:ens:10}Démontrer directement qu'un ensemble à $n$ éléments possède $2^n$ parties.\end{exercice}
\begin{exercice}[Bijection à partir des injections]\label{ex:ens:11}Établir $|(0,1)|=|[0,1]|$ par Cantor--Bernstein.\end{exercice}
\begin{exercice}[Dénombrabilité des mots]\label{ex:ens:12}Montrer que l'ensemble des mots de longueur finie sur un alphabet fini non vide est dénombrable.\end{exercice}
\subsection*{Niveau 4 — Approfondir}
\begin{exercice}[Une diagonale abstraite]\label{ex:ens:13}Pour $F:E\to\PP(E)$, construire explicitement une partie de $E$ qui n'est pas dans l'image de $F$.\end{exercice}
\begin{exercice}[Dénombrabilité des algébriques]\label{ex:ens:14}Montrer que l'ensemble des réels algébriques est dénombrable. En déduire l'existence de réels transcendants.\end{exercice}
\begin{exercice}[Suites binaires]\label{ex:ens:15}Construire une bijection entre $\PP(\N)$ et $\{0,1\}^{\N}$ ; montrer que ce dernier ensemble n'est pas dénombrable.\end{exercice}
\subsection*{Niveau 5 — Défis}
\begin{exercice}[Une bijection inattendue]\label{ex:ens:16}Donner une bijection explicite de $[0,1)$ sur $(0,1)$ en ne déplaçant qu'une suite dénombrable de points.\end{exercice}
\begin{exercice}[Injections et images]\label{ex:ens:17}Montrer qu'une application $f:E\to F$ est injective si et seulement si $f(A\cap B)=f(A)\cap f(B)$ pour toutes parties $A,B\subseteq E$.\end{exercice}
\begin{exercice}[Défi de cardinalité]\label{ex:ens:18}Montrer que $\PP(\N\times\N)$ et $\PP(\N)$ sont équipotents, puis justifier que $\PP(\N)$ n'est pas dénombrable.\end{exercice}


\subsection*{Compléments — modéliser et déjouer les paradoxes}
\begin{exercice}[Graphe, existence et unicité]\label{ex:ens:19}
Soient $A=\{1,2,3\}$ et $B=\{a,b\}$. On considère les relations
\[
 G=\{(1,a),(2,b),(3,a)\},\qquad
 H=\{(1,a),(2,b),(2,a),(3,a)\}.
\]
Sont-elles des graphes d'applications $A\to B$ ? Justifier par les quantificateurs.
\end{exercice}
\begin{exercice}[Le bibliothécaire]\label{ex:ens:20}
Formaliser la règle « le catalogue $C$ recense exactement les catalogues qui ne se recensent pas eux-mêmes ».
En appliquant cette règle à $C$, expliquer le conflit logique.
\end{exercice}
\begin{exercice}[Défi — une relation à réparer]\label{ex:ens:21}
Pour une relation finie $G\subseteq A\times B$ avec $A$ et $B$ non vides, décrire
une procédure qui supprime ou ajoute des couples pour obtenir le graphe d'une application $A\to B$.
Montrer qu'elle termine et expliquer pourquoi l'application obtenue n'est pas nécessairement unique.
\end{exercice}


\begin{retenir}La cardinalité compare des ensembles par bijections ; le calcul propositionnel du chapitre suivant expliquera formellement les règles logiques qui interviennent dans toutes ces preuves.\end{retenir}


\chapter{Calcul propositionnel et calcul des prédicats}
\label{chap:calcul}
\begin{flushleft}\small\sffamily
\textbf{Prérequis :} phrases quantifiées du chapitre 1, ensembles et applications du chapitre 2.\qquad
\textbf{Volume indicatif :} quatre à cinq semaines (cours et TD).
\end{flushleft}
\begin{questionfondatrice}
Peut-on vérifier un raisonnement en ne connaissant que sa forme, sans dépendre du sujet dont parlent les propositions ? Et comment exprimer « tous », « au moins un » et « un seul » sans ambiguïté ?
\end{questionfondatrice}
\begin{objectifs}
Construire une table de vérité, reconnaître tautologie et contradiction, transformer une formule sans changer sa valeur, vérifier la validité d'un argument, manipuler quantificateurs et variables libres, traduire et nier des définitions d'analyse et d'algèbre.
\end{objectifs}

\section{Propositions et connecteurs}
Une lettre propositionnelle $p$ représente une phrase susceptible d'être vraie (V) ou fausse (F). Les connecteurs produisent de nouvelles propositions : négation $\neg p$, conjonction $p\wedge q$, disjonction inclusive $p\vee q$, implication $p\to q$, équivalence $p\leftrightarrow q$.
\begin{definition}{Valeurs de vérité}{valeurs}
L'implication $p\to q$ est \emph{fausse uniquement} quand $p$ est vraie et $q$ fausse. La disjonction est inclusive : $p\vee q$ est vraie si l'un au moins est vrai, éventuellement les deux.
\end{definition}
\begin{center}\renewcommand{\arraystretch}{1.23}
\begin{tabular}{cc|ccccc}\toprule
$p$&$q$&$\neg p$&$p\wedge q$&$p\vee q$&$p\to q$&$p\leftrightarrow q$\\\midrule
V&V&F&V&V&V&V\\
V&F&F&F&V&F&F\\
F&V&V&F&V&V&F\\
F&F&V&F&F&V&V\\\bottomrule
\end{tabular}\end{center}
\begin{exemple}{Pourquoi une hypothèse fausse donne une implication vraie}{vacuite}
« Si $6$ est impair, alors $6^2$ est impair » a une hypothèse fausse ; comme implication matérielle, elle est vraie. Cela n'établit nullement que $6^2$ soit impair. La logique évalue la \emph{forme conditionnelle}, pas la conclusion prise isolément.
\end{exemple}
\begin{figure}[htbp]\centering
\begin{tikzpicture}[>=Latex,node distance=6mm]
\node[draw=LFblue,fill=LFpale,rounded corners=1mm,minimum width=18mm,minimum height=9mm] (p) {$p$};
\node[draw=LFblue,fill=LFpale,rounded corners=1mm,below=11mm of p,minimum width=18mm,minimum height=9mm] (q) {$q$};
\node[draw=LFteal,fill=LFgreen,rounded corners=2mm,right=19mm of p,yshift=-5.5mm,minimum width=23mm,minimum height=16mm] (and) {$\wedge$};
\node[right=16mm of and] (o) {$p\wedge q$};
\draw[->,thick,LFblue] (p.east)-|(and.north west);
\draw[->,thick,LFblue] (q.east)-|(and.south west);
\draw[->,thick,LFteal] (and)--(o);
\end{tikzpicture}
\caption{Un connecteur logique combine les valeurs de ses entrées ; la table de vérité définit son fonctionnement.}
\end{figure}


\subsection{Antisymétrie et implication à antécédent faux}
Une relation $R$ est \emph{antisymétrique} si
\[
 \forall x\,\forall y\;\bigl((xRy\wedge yRx)\Longrightarrow x=y\bigr).
\]
Pour l'ordre strict $<$ sur $\R$, aucun couple $x,y$ ne vérifie simultanément
$x<y$ et $y<x$. L'antécédent est donc toujours faux et l'implication vraie :
la relation $<$ est antisymétrique \emph{au sens logique de cette définition},
mais elle est aussi \emph{irréflexive} et \emph{asymétrique}.
\begin{exemple}{Trois propriétés à ne pas confondre}{strict-antisym}
La relation $\le$ est réflexive, antisymétrique et transitive : c'est un ordre
au sens habituel. La relation $<$ est irréflexive et transitive : c'est un ordre
strict. Elle est également antisymétrique, mais cela ne suffit nullement à en faire
un ordre non strict, car la réflexivité manque.
\end{exemple}
\begin{center}\renewcommand{\arraystretch}{1.2}
\begin{tabular}{lcc}\toprule
Propriété & $\le$ sur $\R$ & $<$ sur $\R$\\\midrule
Réflexive & oui & non\\
Antisymétrique & oui & oui (par vacuité)\\
Asymétrique & non & oui\\
Transitive & oui & oui\\\bottomrule
\end{tabular}\end{center}
\begin{attention}
Antisymétrique ne signifie pas « le contraire de symétrique ».
Une relation asymétrique vérifie $xRy\Rightarrow\neg(yRx)$ ; cela entraîne
l'irréflexivité et l'antisymétrie. La réciproque est fausse.
\end{attention}


\section{Formules, portée et priorité des connecteurs}
Une formule se construit récursivement à partir de lettres $p,q,\ldots$ en appliquant les connecteurs. Nous utilisons l'ordre usuel : $\neg$ avant $\wedge$, puis $\vee$, puis $\to$ ; les parenthèses restent la méthode sûre pour éviter les ambiguïtés.
\begin{definition}{Tautologie, contradiction, équivalence logique}{tautologie}
Une \emph{tautologie} vaut V sous toute attribution des valeurs de ses lettres ; une \emph{contradiction} vaut F partout ; une formule \emph{satisfiable} vaut V dans au moins une attribution. Deux formules $F,G$ sont logiquement équivalentes, noté $F\equiv G$, si $F\leftrightarrow G$ est une tautologie.
\end{definition}
\begin{proposition}{Identités utiles}{identites}
On a
\[
(p\to q)\equiv(\neg p\vee q),\quad
\neg(p\wedge q)\equiv(\neg p\vee\neg q),\quad
\neg(p\vee q)\equiv(\neg p\wedge\neg q).
\]
De plus $p\to q\equiv\neg q\to\neg p$ (contraposée), mais en général $p\to q\not\equiv q\to p$ (réciproque).
\end{proposition}
\begin{proof}
Les formules comparées ont la même valeur dans chacune des quatre lignes du tableau $p,q$. Pour la réciproque, choisir $p=\mathrm F$, $q=\mathrm V$ : $p\to q$ est V, $q\to p$ est F.
\end{proof}
\begin{exemple}{Le piège de l'affirmation du conséquent}{consequent}
« Si $n$ est multiple de $4$, alors $n$ est pair ; $n$ est pair ; donc $n$ est multiple de $4$. » Ce raisonnement a la forme $p\to q,\ q\therefore p$, qui n'est pas valide : $n=2$ le réfute.
\end{exemple}
\begin{attention}
Un énoncé peut être vrai par hasard alors que le schéma d'inférence est invalide. On cherche un contre-modèle pour invalider une \emph{forme} : prémisses toutes vraies et conclusion fausse.
\end{attention}

\section{Règles d'inférence et validité}
\begin{definition}{Argument valide}{argument}
Un argument $P_1,\ldots,P_k\therefore Q$ est \emph{valide} lorsqu'il n'existe aucune attribution rendant vraies toutes les prémisses et fausse la conclusion. Cela équivaut à dire que $(P_1\wedge\cdots\wedge P_k)\to Q$ est une tautologie.
\end{definition}
\begin{center}\renewcommand{\arraystretch}{1.3}
\begin{tabularx}{\textwidth}{>{\bfseries}l X X}\toprule
Règle&Prémisses&Conclusion\\\midrule
Modus ponens&$p\to q$, $p$&$q$\\
Modus tollens&$p\to q$, $\neg q$&$\neg p$\\
Syllogisme hypothétique&$p\to q$, $q\to r$&$p\to r$\\
Élimination de $\wedge$&$p\wedge q$&$p$\\
Introduction de $\vee$&$p$&$p\vee q$\\\bottomrule
\end{tabularx}\end{center}
\begin{exemple}{Une preuve en trois pas}{trois-pas}
Prémisses : $p\to q$, $q\to r$, $p$. Par modus ponens on obtient $q$ ; puis avec $q\to r$, on obtient $r$. Il ne suffit pas d'affirmer que la conclusion « semble logique » : chaque pas possède une règle.
\end{exemple}
\begin{proposition}{Raisonnement par disjonction de cas}{cas-disj}
Les prémisses $p\vee q$, $p\to r$ et $q\to r$ entraînent $r$.
\end{proposition}
\begin{proof}Si $p\vee q$ est vraie, au moins une branche est vraie. Dans la première, $p\to r$ donne $r$ ; dans la seconde, $q\to r$ donne encore $r$.\end{proof}

\section{Formes normales et méthode systématique}
Une \emph{forme normale disjonctive} (FND) est une disjonction de conjonctions de littéraux ; une \emph{forme normale conjonctive} (FNC) est une conjonction de disjonctions de littéraux. On peut construire une FND à partir des lignes vraies de la table de vérité, et une FNC à partir des lignes fausses.
\begin{exemple}{Construire une FND}{fnd}
La formule $p\leftrightarrow q$ est vraie aux lignes $(\mathrm V,\mathrm V)$ et $(\mathrm F,\mathrm F)$, d'où
\[(p\wedge q)\vee(\neg p\wedge\neg q).\]
La même formule possède la FNC $(\neg p\vee q)\wedge(\neg q\vee p)$.
\end{exemple}
\begin{exemple}{Tester une argumentation avec un contre-modèle}{contre-modele}
Pour $p\to q,\ q\therefore p$, l'attribution $p=\mathrm F$, $q=\mathrm V$ rend les deux prémisses vraies et la conclusion fausse. Une seule ligne suffit à invalider l'argument.
\end{exemple}

\section{Prédicats et quantificateurs}
Une expression $P(x)$ est une \emph{formule ouverte} si $x$ est libre. Elle devient proposition lorsque l'on fixe $x$ ou qu'on quantifie cette variable dans un domaine $E$. La formule $\forall x\in E\ P(x)$ affirme que tous satisfont $P$ ; $\exists x\in E\ P(x)$ qu'au moins un le satisfait.
\begin{definition}{Variables libres et liées}{liees}
Une occurrence de variable est \emph{liée} lorsqu'elle se trouve dans la portée d'un quantificateur qui la déclare. Sinon elle est \emph{libre}. Une formule sans variables libres est dite fermée. La portée se lit à l'aide des parenthèses.
\end{definition}
\begin{exemple}{Changer l'ordre change la phrase}{ordre-quant}
Sur $\R$, $\forall x\,\exists y\ (y>x)$ est vraie, en choisissant $y=x+1$. La phrase $\exists y\,\forall x\ (y>x)$ est fausse : il n'existe pas de majorant universel de $\R$.
\end{exemple}
\begin{figure}[htbp]\centering
\begin{tikzpicture}[x=.85cm,y=.82cm]
\foreach \x in {0,...,4}{\foreach \y in {0,...,3}{\draw[LFblue!40] (\x,\y) rectangle +(1,1);}}
\foreach \x in {0,...,4}{\node[below,font=\footnotesize] at (\x+.5,-.12) {$x_{\the\numexpr\x+1\relax}$};}
\foreach \y in {0,...,3}{\node[left,font=\footnotesize] at (-.05,\y+.5) {$y_{\the\numexpr\y+1\relax}$};}
\foreach \x/\y in {0/0,1/2,2/1,3/3,4/0}{\fill[LFteal!55] (\x+.04,\y+.04) rectangle (\x+.96,\y+.96);}
\node[align=left,anchor=west,font=\small] at (5.45,2.8) {$\forall x\,\exists y\,P(x,y)$ :};
\node[align=left,anchor=west,font=\small] at (5.45,2.15) {chaque colonne contient};
\node[align=left,anchor=west,font=\small] at (5.45,1.65) {au moins une case colorée.};
\node[align=left,anchor=west,font=\small] at (5.45,.8) {$\exists y\,\forall x\,P(x,y)$ :};
\node[align=left,anchor=west,font=\small] at (5.45,.15) {une ligne entière colorée.};
\end{tikzpicture}
\caption{Les quantificateurs échangés imposent des configurations différentes.}
\end{figure}
\begin{proposition}{Négations quantifiées}{neg-quant}
\[
\neg\forall x\,P(x)\equiv\exists x\,\neg P(x),\qquad
\neg\exists x\,P(x)\equiv\forall x\,\neg P(x).
\]
Le domaine de quantification est conservé. Pour nier une chaîne, on inverse les quantificateurs un par un et l'on nie la propriété finale.
\end{proposition}
\begin{proof}La première négation signifie qu'au moins un élément contredit $P$ ; la seconde qu'aucun élément ne satisfait $P$.\end{proof}
\begin{exemple}{Nier la continuité en un point}{contin-niee}
Pour $f:\R\to\R$, la continuité en $a$ est
\[\forall\varepsilon>0\ \exists\delta>0\ \forall x\in\R:\ |x-a|<\delta\ \Longrightarrow\ |f(x)-f(a)|<\varepsilon.\]
Sa négation est
\[\exists\varepsilon_0>0\ \forall\delta>0\ \exists x\in\R:\ |x-a|<\delta\ \wedge\ |f(x)-f(a)|\ge\varepsilon_0.\]
Le « et » final provient de $\neg(P\to Q)\equiv P\wedge\neg Q$.
\end{exemple}
\begin{definition}{Existence et unicité}{unique}
$\exists!x\,P(x)$ signifie $\exists x\,[P(x)\wedge\forall y(P(y)\to y=x)]$. Pour démontrer l'unicité, on suppose $P(x)$ et $P(y)$ et on prouve $x=y$.
\end{definition}
\begin{attention}
La substitution d'une expression contenant une variable dans une formule quantifiée doit éviter la \emph{capture} de variable. Il est permis de renommer une variable liée : $\forall x\,P(x)$ et $\forall z\,P(z)$ expriment la même chose, si le remplacement est cohérent.
\end{attention}

\section{Exercices progressifs}
\subsection*{Niveau 1 — Lire les formules}
\begin{exercice}[Table de vérité]\label{ex:prop:1}Construire la table de $p\to(q\vee\neg p)$. Est-ce une tautologie ?\end{exercice}
\begin{exercice}[Négations]\label{ex:prop:2}Simplifier $\neg(p\to q)$ et $\neg(p\vee\neg q)$.\end{exercice}
\begin{exercice}[Quantificateurs]\label{ex:prop:3}Évaluer dans $\Z$ : $\forall n\,\exists m\ (m>n)$ et $\exists m\,\forall n\ (m>n)$.\end{exercice}
\begin{exercice}[Variables libres]\label{ex:prop:4}Identifier les variables libres de $\forall x\,(P(x,y)\to\exists z\,Q(z,x))$.\end{exercice}
\subsection*{Niveau 2 — Calculer et prouver}
\begin{exercice}[Contraposée]\label{ex:prop:5}Vérifier par tableau $p\to q\equiv\neg q\to\neg p$.\end{exercice}
\begin{exercice}[Équivalence]\label{ex:prop:6}Démontrer $p\to q\equiv\neg p\vee q$ et donner sa FNC.\end{exercice}
\begin{exercice}[Inférence]\label{ex:prop:7}À partir de $p\to q$, $q\to r$, $p$, déduire $r$ en nommant les règles.\end{exercice}
\begin{exercice}[Unicité]\label{ex:prop:8}Traduire « $x^2=0$ admet une unique solution réelle » avec $\exists!$, puis sans ce symbole.\end{exercice}
\subsection*{Niveau 3 — Formaliser}
\begin{exercice}[Argument invalide]\label{ex:prop:9}Trouver un contre-modèle pour $p\to q,\ q\therefore p$.\end{exercice}
\begin{exercice}[Forme normale]\label{ex:prop:10}Trouver une FND et une FNC de $p\leftrightarrow q$.\end{exercice}
\begin{exercice}[Ordre des quantificateurs]\label{ex:prop:11}Sur $\N$, interpréter $\forall x\exists y\ (x<y)$ et $\exists y\forall x\ (x<y)$ ; justifier leurs valeurs de vérité.\end{exercice}
\begin{exercice}[Négation d'une limite]\label{ex:prop:12}Écrire formellement puis nier « $u_n\to\ell$ » (suite réelle).\end{exercice}
\subsection*{Niveau 4 — Approfondir}
\begin{exercice}[Règle par cas]\label{ex:prop:13}Prouver que $(p\vee q)\wedge(p\to r)\wedge(q\to r)$ entraîne $r$.\end{exercice}
\begin{exercice}[Validité et satisfiabilité]\label{ex:prop:14}Les formules $p\vee q$, $\neg p$ et $\neg q$ sont-elles simultanément satisfiables ? Justifier.\end{exercice}
\begin{exercice}[Négation de l'uniforme]\label{ex:prop:15}Écrire et nier « $f_n\to f$ uniformément sur $E$ », avec $f_n,f:E\to\R$.\end{exercice}
\subsection*{Niveau 5 — Défis}
\begin{exercice}[Exactement un des trois]\label{ex:prop:16}Écrire une formule en $p,q,r$ vraie exactement lorsque l'une des trois lettres est vraie ; simplifier sa FND.\end{exercice}
\begin{exercice}[Satisfaction sans calcul exhaustif]\label{ex:prop:17}La formule $(p\to q)\wedge(q\to r)\wedge(r\to p)\wedge p\wedge\neg r$ est-elle satisfiable ? Démontrer sans énumérer les huit lignes.\end{exercice}
\begin{exercice}[Deux propositions très proches]\label{ex:prop:18}Dans $\R$, comparer $\forall x\exists y\ (x+y=0)$ et $\exists y\forall x\ (x+y=0)$. Formaliser leur négation respective et établir leur valeur de vérité.\end{exercice}


\subsection*{Compléments — relations et implication}
\begin{exercice}[L'ordre strict]\label{ex:prop:19}
Écrire avec quantificateurs l'antisymétrie de $<$ sur $\R$ et justifier sa valeur de vérité.
La relation est-elle réflexive ? Est-elle un ordre au sens non strict ?
\end{exercice}
\begin{exercice}[Asymétrie versus antisymétrie]\label{ex:prop:20}
Sur $E=\{a,b\}$, donner une relation antisymétrique qui ne soit pas asymétrique.
\end{exercice}
\begin{exercice}[Défi — les quatre combinaisons]\label{ex:prop:21}
Trouver sur un ensemble fini un exemple de relation (i) symétrique et antisymétrique,
(ii) symétrique mais non antisymétrique, (iii) antisymétrique mais non symétrique,
(iv) ni symétrique ni antisymétrique. Justifier chaque réponse.
\end{exercice}


\begin{retenir}Tables de vérité et raisonnements quantifiés formalisent les preuves du chapitre 1. Le chapitre 4 emploiera ces règles pour démontrer des propriétés valables pour une infinité d'entiers.\end{retenir}

\chapter{Bon ordre et preuve par récurrence}
\label{chap:recurrence}
\begin{flushleft}\small\sffamily
\textbf{Prérequis :} implications, ensembles ordonnés, quantificateurs.\qquad
\textbf{Volume indicatif :} trois à quatre semaines (cours et TD).
\end{flushleft}
\begin{questionfondatrice}
Comment une vérification initiale et une règle permettant de passer d'un entier au suivant peuvent-elles démontrer une infinité de cas ? Que faire lorsqu'un contre-exemple peut être choisi minimal ?
\end{questionfondatrice}
\begin{objectifs}
Savoir choisir la bonne propriété de récurrence ; distinguer récurrence simple, forte et simultanée ; exploiter le principe du plus petit élément ; reconnaître ordre partiel et bon ordre ; comprendre le statut du théorème du bon ordre de Zermelo.
\end{objectifs}

\section{Le principe de récurrence}
\begin{theoreme}{Récurrence simple}{simple}
Soit $P(n)$ une propriété pour $n\in\N$. Si $P(0)$ est vraie et si, pour tout $n\in\N$, $P(n)\Rightarrow P(n+1)$, alors $P(n)$ est vraie pour tout $n\in\N$.
\end{theoreme}
\begin{proof}
Supposons que l'ensemble $C=\{n\in\N:\neg P(n)\}$ soit non vide, et prenons son plus petit élément $m$. Comme $P(0)$, $m\ne0$, donc $m-1\in\N$ et $P(m-1)$ est vraie par minimalité. L'hérédité donne $P(m)$ : contradiction. Cette preuve emploie le principe du bon ordre de $\N$, étudié plus loin.
\end{proof}
\begin{figure}[htbp]\centering
\begin{tikzpicture}[>=Latex]
\foreach \i in {0,...,5}{
 \draw[LFblue,thick,fill=LFpale,rounded corners=1mm] (1.18*\i,0) rectangle +(0.82,1.1);
 \node[font=\small] at (1.18*\i+.41,.55) {$P(\i)$};
}
\foreach \i in {0,...,4}{\draw[->,LFteal,thick] (1.18*\i+.84,.55)--(1.18*\i+1.13,.55);}
\node[above,font=\small,LFteal] at (.41,1.2) {initialisation};
\node[above,font=\small,LFteal] at (3.3,1.2) {hérédité};
\node[right,LFnavy] at (7.05,.55) {$\cdots$};
\end{tikzpicture}
\caption{Le premier domino ne suffit pas : il faut aussi un mécanisme valable pour chaque passage $n\to n+1$.}
\end{figure}
\begin{exemple}{La somme des premiers entiers}{somme-entiers}
Pour $n\ge0$, démontrons $\sum_{k=0}^n k=\frac{n(n+1)}2$. À $n=0$, les deux membres valent $0$. Si l'identité est vraie au rang $n$, alors
\[
\sum_{k=0}^{n+1}k=\frac{n(n+1)}2+(n+1)=\frac{(n+1)(n+2)}2.
\]
La propriété est donc vraie pour tous les $n$.
\end{exemple}
\begin{attention}
L'hérédité ne remplace pas l'initialisation. La propriété $P(n)$ : « $n<0$ » vérifie formellement $P(n)\Rightarrow P(n+1)$ pour tout $n\in\N$ (hypothèse toujours fausse), mais elle n'est vraie à aucun rang.
\end{attention}

\section{Récurrence forte et récurrence à plusieurs rangs}
\begin{theoreme}{Récurrence forte}{forte}
Si $P(0)$ est vraie et si, pour chaque $n\in\N$, les hypothèses $P(0),\ldots,P(n)$ entraînent $P(n+1)$, alors $P(n)$ vaut pour tout $n$.
\end{theoreme}
\begin{proof}
Posons $Q(n)=\bigwedge_{k=0}^nP(k)$. Alors $Q(0)$ est vraie et $Q(n)\Rightarrow Q(n+1)$ par l'hypothèse forte. Appliquer la récurrence simple à $Q$.
\end{proof}
\begin{exemple}{Chaque entier se factorise en nombres premiers}{factorisation}
Pour tout $n\ge2$, soit $n$ premier, soit $n=ab$ avec $2\le a,b<n$. En supposant les entiers de $2$ à $n-1$ déjà décomposables en produit de premiers, il en va de même de $a,b$, donc de $n$. Le cas initial $n=2$ est premier.
\end{exemple}
Une récurrence d'ordre deux nécessite deux initialisations : pour $P(n+2)$ à partir de $P(n)$ et $P(n+1)$, établir $P(0)$ et $P(1)$. La récurrence simultanée consiste à grouper deux propriétés sous $Q(n)=P(n)\wedge R(n)$.
\begin{exemple}{Suite de Fibonacci}{fibonacci}
Si $F_0=0$, $F_1=1$, $F_{n+2}=F_{n+1}+F_n$, alors $F_n\le2^{n-1}$ pour $n\ge1$. Les rangs $1$ et $2$ vérifient l'inégalité ; si elle est vraie aux deux rangs précédents, alors $F_{n+2}\le2^n+2^{n-1}\le2^{n+1}$. On peut aussi obtenir la majoration par une récurrence forte.
\end{exemple}


\section{Approfondissement : récurrence de Cauchy}
\begin{theoreme}{Principe de récurrence de Cauchy}{induction-cauchy}
Soit $P(n)$ une propriété définie pour tout entier $n\ge1$. Supposons que :
\begin{enumerate}
 \item $P(1)$ est vraie ;
 \item pour tout $n\ge1$, $P(n)\Rightarrow P(2n)$ ;
 \item pour tout $n\ge2$, $P(n)\Rightarrow P(n-1)$.
\end{enumerate}
Alors $P(n)$ est vraie pour tout $n\ge1$.
\end{theoreme}
\begin{proof}
Par doublement successif, $P(1),P(2),P(4),\ldots,P(2^k)$ sont vraies.
Fixons $m\ge1$ et choisissons $k$ tel que $2^k\ge m$. En utilisant
la troisième hypothèse un nombre fini de fois à partir de $P(2^k)$,
on obtient $P(2^k-1),\ldots,P(m)$. Comme $m$ est arbitraire, la preuve est achevée.
\end{proof}
\begin{figure}[htbp]\centering
\begin{tikzpicture}[>=Latex,node distance=9mm]
\foreach \x/\s in {0/1,1.7/2,3.4/4,5.1/8}{
 \node[draw=LFblue,fill=LFpale,rounded corners=2mm,minimum width=12mm,minimum height=7mm] (n\s) at (\x,0) {$P(\s)$};}
\draw[->,LFteal,thick] (n1)--(n2); \draw[->,LFteal,thick] (n2)--(n4);\draw[->,LFteal,thick] (n4)--(n8);
\node[draw=LFgold,fill=LFpaper,rounded corners=2mm] (seven) at (5.1,-1.4) {$P(7)$};
\node[draw=LFgold,fill=LFpaper,rounded corners=2mm] (six) at (3.4,-1.4) {$P(6)$};
\node[draw=LFgold,fill=LFpaper,rounded corners=2mm] (five) at (1.7,-1.4) {$P(5)$};
\draw[->,LFgold,thick] (n8)--(seven);\draw[->,LFgold,thick] (seven)--(six);\draw[->,LFgold,thick] (six)--(five);
\node[above=2mm of n4,font=\small,LFteal]{doublement};
\node[below=2mm of six,font=\small,LFgold!70!black]{descente};
\end{tikzpicture}
\caption{Deux chemins de propagation : on atteint une puissance de $2$, puis on redescend au rang voulu.}
\end{figure}
\begin{exemple}{Inégalité arithmético-géométrique par récurrence de Cauchy}{agm-cauchy}
Pour des nombres $a_1,\ldots,a_n>0$, la propriété $P(n)$ est
\[
 (a_1\cdots a_n)^{1/n}\le\frac{a_1+\cdots+a_n}{n}.
\]
$P(1)$ est une égalité. Pour passer de $n$ à $2n$, appliquer $P(n)$ aux deux
moitiés, de moyennes respectives $A$ et $B$. Le produit total est au plus
$(AB)^n$, puis $AB\le((A+B)/2)^2$ montre $P(2n)$.
Pour redescendre de $n$ à $n-1$, partir de $a_1,\ldots,a_{n-1}$ de moyenne
$m>0$ et leur ajouter $a_n=m$. L'hypothèse $P(n)$ donne
$(a_1\cdots a_{n-1}m)^{1/n}\le m$, donc
$a_1\cdots a_{n-1}\le m^{n-1}$ : c'est $P(n-1)$.
Le principe de Cauchy conclut à tous les rangs.
\end{exemple}

\section{Approfondissement : Cauchy--Schwarz par récurrence}
\begin{theoreme}{Inégalité de Cauchy--Schwarz finie}{cauchy-schwarz-ind}
Pour tous $(a_i)_{1\le i\le n},(b_i)_{1\le i\le n}\in\R^n$,
\[
 \left(\sum_{i=1}^n a_ib_i\right)^2\le
 \left(\sum_{i=1}^n a_i^2\right)
 \left(\sum_{i=1}^n b_i^2\right).
\]
\end{theoreme}
\begin{proof}
Pour $n=1$, il y a égalité. Supposons l'inégalité vraie au rang $n$
et posons $A=\sum_{i=1}^n a_i^2$, $B=\sum_{i=1}^n b_i^2$ et
$S=\sum_{i=1}^n a_ib_i$. En écrivant $a=a_{n+1}$ et $b=b_{n+1}$,
\[
\begin{aligned}
 (A+a^2)(B+b^2)-(S+ab)^2
 &=\underbrace{(AB-S^2)}_{\ge0\text{ par hypothèse}}+
   \underbrace{(Ab^2+Ba^2-2abS)}_{=\sum_{i=1}^n(a_i b-a b_i)^2\ge0}.
\end{aligned}
\]
Les deux termes sont positifs, d'où l'inégalité au rang $n+1$.
\end{proof}
\begin{exemple}{Une application numérique}{cs-num}
Pour $a=(1,2,-1)$ et $b=(2,0,3)$, on trouve $a\cdot b=-1$,
$\sum a_i^2=6$ et $\sum b_i^2=13$. Ainsi $1\le 78$, conformément à l'inégalité.
\end{exemple}
\begin{retenir}
La récurrence de Cauchy est une méthode particulière de propagation des indices,
non une nouvelle inégalité. La preuve de Cauchy--Schwarz ci-dessus utilise, quant à elle,
la récurrence simple et une identité de sommes de carrés.
\end{retenir}


\section{Relations d'ordre, ordres partiels et bons ordres}
\begin{definition}{Ordre et ordre total}{ordre}
Une relation $\preceq$ sur $E$ est un \emph{ordre} si elle est réflexive, antisymétrique et transitive. Elle est \emph{totale} si pour tous $x,y$, $x\preceq y$ ou $y\preceq x$. Un élément $m$ est \emph{minimal} s'il n'existe aucun $x\ne m$ avec $x\preceq m$ ; il est \emph{minimum} si $m\preceq x$ pour tout $x$.
\end{definition}
\begin{exemple}{Deux éléments incomparables}{divisibilite}
Sur les diviseurs positifs de $12$, la divisibilité $a\preceq b\iff a\mid b$ est un ordre partiel : $2$ et $3$ sont incomparables. Le minimum est $1$, le maximum est $12$.
\end{exemple}
\begin{figure}[htbp]\centering
\begin{tikzpicture}[>=Latex,every node/.style={circle,draw=LFblue,fill=LFpale,minimum size=8mm,font=\small}]
\node (one) at (0,0) {$1$};
\node (two) at (-1.8,1.2) {$2$};\node (three) at (1.8,1.2) {$3$};
\node (four) at (-1.8,2.4) {$4$};\node (six) at (1.1,2.4) {$6$};
\node (twelve) at (0,3.7) {$12$};
\foreach \a/\b in {one/two,one/three,two/four,two/six,three/six,four/twelve,six/twelve}{\draw[LFteal,thick] (\a)--(\b);}
\end{tikzpicture}
\caption{Diagramme de Hasse de la divisibilité sur $\{1,2,3,4,6,12\}$. Une arête encode une relation de couverture ; les relations transitives sont implicites.}
\end{figure}
\begin{definition}{Bon ordre}{bon-ordre}
Un ordre total $\preceq$ sur $E$ est un \emph{bon ordre} si chaque partie non vide de $E$ possède un plus petit élément. L'ordre usuel sur $\N$ en est un ; celui sur $\Z$ ne l'est pas, car $\Z$ n'a pas de plus petit élément.
\end{definition}
\begin{attention}
« Tout sous-ensemble non vide a un élément minimal » ne suffit pas à garantir un bon ordre si l'ordre est seulement partiel : sur un ensemble fini de deux éléments incomparables, les deux sont minimaux, mais aucun n'est minimum.
\end{attention}

\section{Le plus petit contre-exemple et les preuves arithmétiques}
\begin{proposition}{Récurrence et plus petit élément}{equivalence-induction}
Sur $\N$, le principe de récurrence et le principe du plus petit élément sont équivalents.
\end{proposition}
\begin{proof}
Le sens « plus petit élément $\Rightarrow$ récurrence » est la preuve du premier théorème. Inversement, supposons la récurrence et prenons une partie non vide $A\subseteq\N$ sans minimum. Posons $P(n)$ : « $A\cap\{0,\ldots,n\}=\varnothing$ ». Comme $0$ ne peut appartenir à $A$ sans être minimum, $P(0)$ vaut. Si $P(n)$ et $n+1\in A$, alors $n+1$ serait le minimum de $A$ ; impossible. Donc $P(n+1)$. La récurrence conduit à $A=\varnothing$, contradiction.
\end{proof}
\begin{exemple}{Infinité des nombres premiers}{infinitude-primes}
Supposons qu'il n'existe qu'un nombre fini de premiers $p_1,\ldots,p_r$. L'entier $N=p_1\cdots p_r+1$ est supérieur à $1$ et admet donc un diviseur premier $p_i$. Ce dernier diviserait à la fois $N$ et le produit $p_1\cdots p_r$, donc leur différence $1$ : contradiction.
\end{exemple}
\begin{theoreme}{Division euclidienne dans $\N$}{division}
Pour $a\in\N$ et $b\in\N^*$, il existe un unique couple $(q,r)\in\N^2$ tel que $a=bq+r$ et $0\le r<b$.
\end{theoreme}
\begin{proof}
L'ensemble $S=\{a-bk:k\in\N,\ a-bk\ge0\}$ est non vide (prendre $k=0$). Soit $r$ son minimum et $q$ tel que $r=a-bq$. Si $r\ge b$, alors $r-b=a-b(q+1)\in S$, strictement plus petit : impossible. Pour l'unicité, si $bq+r=bq'+r'$, alors $b(q-q')=r'-r$ et $|r'-r|<b$ ; le seul multiple de $b$ possible est $0$, d'où $q=q'$ et $r=r'$.
\end{proof}
\begin{theoreme}{Unicité de la décomposition en facteurs premiers}{facteurs-unique}
Tout entier $n\ge2$ s'écrit comme produit de nombres premiers ; cette décomposition est unique à l'ordre des facteurs près.
\end{theoreme}
\begin{proof}
L'existence vient de la récurrence forte. Pour l'unicité, on utilise le lemme d'Euclide : si $p$ est premier et $p\mid ab$, alors $p\mid a$ ou $p\mid b$. Ce lemme découle de Bézout : si $p\nmid a$, alors $\gcd(p,a)=1$ et il existe $u,v\in\Z$ tels que $up+va=1$ ; en multipliant par $b$, on obtient $p\mid b$. Le théorème de Bézout se démontre à partir de la division euclidienne et de l'algorithme d'Euclide. Comparons deux factorisations $p_1\cdots p_r=q_1\cdots q_s$ : $p_1$ divise un des $q_j$ ; étant premiers, ils sont égaux. On les simplifie et on recommence, ce qui établit l'unicité.
\end{proof}

\section{Bon ordre de Zermelo et axiome du choix}
\begin{theoreme}{Théorème du bon ordre de Zermelo (énoncé)}{zermelo}
Dans la théorie ZF, l'axiome du choix équivaut à l'assertion : tout ensemble peut être muni d'un bon ordre. Il équivaut aussi au lemme de Zorn.
\end{theoreme}
\begin{exemple}{Pourquoi l'ordre habituel ne suffit pas}{reels-ordre}
L'ordre usuel sur $[0,1]$ n'est pas un bon ordre, puisque $(0,1]$ n'a pas de plus petit élément. L'affirmation de Zermelo garantit l'existence d'\emph{un autre ordre} sur cet ensemble si l'on admet AC ; elle ne dit pas que l'ordre usuel convient.
\end{exemple}
La preuve générale du théorème de Zermelo requiert une étude plus avancée des ordinaux ou des arguments équivalents. En L2, l'objectif est de connaître un énoncé rigoureux, son hypothèse axiomatique et quelques usages. Il faut distinguer ce résultat du bon ordre naturel de $\N$, utilisé dans nos récurrences, qui ne nécessite pas l'axiome du choix pour son établissement usuel.
\begin{retenir}
Récurrence simple, récurrence forte et méthode du plus petit contre-exemple constituent trois présentations d'une même structure sur $\N$. Le théorème de Zermelo est un résultat d'une portée bien plus large, conditionné par l'axiome du choix dans ZF.
\end{retenir}

\section{Exercices progressifs}
\subsection*{Niveau 1 — Initialiser et transmettre}
\begin{exercice}[Somme d'entiers]\label{ex:rec:1}Montrer $1+\cdots+n=n(n+1)/2$ pour $n\ge1$.\end{exercice}
\begin{exercice}[Somme géométrique]\label{ex:rec:2}Montrer $1+2+\cdots+2^n=2^{n+1}-1$.\end{exercice}
\begin{exercice}[Parité]\label{ex:rec:3}Montrer que $n(n+1)$ est pair pour tout $n\in\N$.\end{exercice}
\begin{exercice}[Trouver le défaut]\label{ex:rec:4}Pourquoi la phrase « $P(n)\Rightarrow P(n+1)$, donc $P(n)$ pour tout $n$ » est-elle incomplète ?\end{exercice}
\subsection*{Niveau 2 — Construire une preuve}
\begin{exercice}[Somme des carrés]\label{ex:rec:5}Prouver $\sum_{k=1}^n k^2=n(n+1)(2n+1)/6$.\end{exercice}
\begin{exercice}[Divisibilité]\label{ex:rec:6}Montrer que $7\mid(8^n-1)$ pour tout $n\ge0$.\end{exercice}
\begin{exercice}[Ordre partiel]\label{ex:rec:7}Sur $\PP(E)$, vérifier que $\subseteq$ est un ordre et dire quand il est total.\end{exercice}
\begin{exercice}[Minimum ou minimal]\label{ex:rec:8}Sur les diviseurs positifs de $12$ ordonnés par divisibilité, déterminer minimum, maximum et une paire incomparable.\end{exercice}
\subsection*{Niveau 3 — Choisir une méthode}
\begin{exercice}[Récurrence forte]\label{ex:rec:9}Montrer que tout $n\ge2$ possède un diviseur premier.\end{exercice}
\begin{exercice}[Plus petit contre-exemple]\label{ex:rec:10}À partir du bon ordre de $\N$, redémontrer la récurrence simple.\end{exercice}
\begin{exercice}[Fibonacci]\label{ex:rec:11}Avec $F_0=0$, $F_1=1$, $F_{n+2}=F_{n+1}+F_n$, montrer $F_n\le2^{n-1}$ pour tout $n\ge1$.\end{exercice}
\begin{exercice}[Un bon ordre modifié]\label{ex:rec:12}L'ordre inverse sur $\N$ est-il un bon ordre ? Justifier.\end{exercice}
\subsection*{Niveau 4 — Approfondir}
\begin{exercice}[Division euclidienne]\label{ex:rec:13}À l'aide du plus petit élément, démontrer l'existence du reste de la division de $a\in\N$ par $b\ge1$.\end{exercice}
\begin{exercice}[Somme pondérée]\label{ex:rec:14}Prouver $\sum_{k=0}^{n}k2^k=(n-1)2^{n+1}+2$ pour $n\ge0$.\end{exercice}
\begin{exercice}[Inégalité factorielle]\label{ex:rec:15}Montrer $2^n<n!$ pour tout $n\ge4$.\end{exercice}
\subsection*{Niveau 5 — Défis}
\begin{exercice}[Écriture binaire]\label{ex:rec:16}Prouver par récurrence forte que tout entier $n\ge1$ admet une écriture binaire finie, puis en établir l'unicité.\end{exercice}
\begin{exercice}[Algorithme d'Euclide]\label{ex:rec:17}Montrer que l'algorithme des restes successifs termine et que son dernier reste non nul est le PGCD des deux entiers initiaux.\end{exercice}
\begin{exercice}[Un invariant de récurrence]\label{ex:rec:18}Soit $u_0=1$ et $u_{n+1}=2u_n+1$. Deviner et prouver une expression explicite de $u_n$. Expliquer pourquoi $u_n$ est impair pour tout $n$.\end{exercice}


\subsection*{Compléments — Cauchy et Cauchy--Schwarz}
\begin{exercice}[Récurrence de Cauchy]\label{ex:rec:19}
Supposons $P(1)$, $P(n)\Rightarrow P(2n)$ pour $n\ge1$ et $P(n)\Rightarrow P(n-1)$ pour $n\ge2$.
Décrire une chaîne de déductions qui établit $P(13)$.
\end{exercice}
\begin{exercice}[Cauchy--Schwarz en dimension trois]\label{ex:rec:20}
Vérifier l'inégalité de Cauchy--Schwarz pour $a=(1,0,2)$ et $b=(2,1,-1)$.
\end{exercice}
\begin{exercice}[Défi — le cas d'égalité]\label{ex:rec:21}
À partir de l'identité utilisée dans la preuve par récurrence de Cauchy--Schwarz,
montrer que pour deux vecteurs réels, l'égalité a lieu si et seulement si les vecteurs
sont linéairement dépendants. Traiter aussi le cas où l'un des vecteurs est nul.
\end{exercice}


\begin{retenir}Une preuve par récurrence est une preuve complète, non une extrapolation empirique. Il faut vérifier les cas initiaux, énoncer précisément l'hypothèse de récurrence et justifier le passage au rang suivant.\end{retenir}

\backmatter
\chapter*{Bibliographie générale}
\addcontentsline{toc}{chapter}{Bibliographie générale}
\markboth{Bibliographie générale}{Bibliographie générale}
Les références ci-dessous sont proposées comme ouvrages de travail et ressources d'approfondissement pour l'ensemble du cours. Les exercices du présent polycopié sont reformulés à partir de ces sources.
\begin{itemize}[leftmargin=0pt,label={},itemsep=9pt]
\item R. Cori et D. Lascar, \emph{Logique mathématique, tome 1 : Calcul propositionnel, algèbre de Boole, calcul des prédicats}, Dunod, Paris, 2003.
\item S. Epp, \emph{Discrete Mathematics with Applications}, 4th ed., Brooks/Cole, 2011.
\item H. B. Enderton, \emph{Elements of Set Theory}, Academic Press, New York, 1977.
\item P. R. Halmos, \emph{Naive Set Theory}, D. Van Nostrand, Princeton, 1960.
\item R. Hammack, \emph{Book of Proof}, 3rd ed., 2018 (version révisée en 2025). Version PDF libre : \url{https://richardhammack.github.io/BookOfProof/}.
\item J.-L. Krivine, \emph{Théorie des ensembles}, Cassini, Paris, 1998.
\item E. Lehman, F. T. Leighton et A. R. Meyer, \emph{Mathematics for Computer Science}, MIT OpenCourseWare, cours en ligne : \href{https://ocw.mit.edu/courses/6-1200j-mathematics-for-computer-science-spring-2024/}{MIT OpenCourseWare}.
\item K. H. Rosen, \emph{Discrete Mathematics and Its Applications}, 8th ed., McGraw-Hill, 2019.
\item D. J. Velleman, \emph{How to Prove It: A Structured Approach}, 3rd ed., Cambridge University Press, 2019.
\item R. Chemlal, \emph{Cours de logique mathématique}, polycopié de Licence 2, Université de Béjaïa, version institutionnelle : \href{https://elearning.univ-bejaia.dz/pluginfile.php/503483/mod_resource/content/0/Cours_CHEMLAL%20Rezki_Logique%20Math%C3%A9matique.pdf}{Université de Béjaïa}.
\end{itemize}

\end{document}
